\documentclass[11pt,a4paper]{article}

\usepackage[margin=2.0cm]{geometry}
\usepackage{amsmath,amssymb,amsthm,amsfonts,mathtools}
\usepackage[T1]{fontenc}
\usepackage{lmodern}
\usepackage{microtype}
\usepackage{enumitem}
\usepackage{booktabs}
\usepackage{xcolor}
\usepackage[hidelinks,colorlinks=true,
            urlcolor=blue!50!black,linkcolor=blue!50!black]{hyperref}
\usepackage{listings}
\usepackage{tikz}
\usepackage{tikz-cd}
\usetikzlibrary{positioning,arrows.meta,shapes.geometric,fit,calc,
                decorations.pathreplacing,backgrounds}

\definecolor{matclr}{RGB}{40,103,160}
\definecolor{unitclr}{RGB}{180,70,40}
\definecolor{rawclr}{RGB}{50,140,80}
\definecolor{prodclr}{RGB}{160,60,140}
\definecolor{kept}{RGB}{40,140,80}
\definecolor{cut}{RGB}{170,60,60}
\definecolor{accent}{RGB}{200,120,30}

\theoremstyle{plain}
\newtheorem{theorem}{Theorem}
\newtheorem{proposition}[theorem]{Proposition}
\newtheorem{lemma}[theorem]{Lemma}
\newtheorem{corollary}[theorem]{Corollary}

\theoremstyle{definition}
\newtheorem{definition}[theorem]{Definition}
\newtheorem{remark}[theorem]{Remark}
\newtheorem{example}[theorem]{Example}
\newtheorem{axiom}{Axiom}[section]
\renewcommand{\theaxiom}{A\arabic{axiom}}

\lstset{
  basicstyle=\ttfamily\footnotesize,
  breaklines=true, columns=fullflexible, keepspaces=true,
  showstringspaces=false, frame=single, rulecolor=\color{gray!30},
}

\newcommand{\HyMeKo}{\textsc{HyMeKo}}
\newcommand{\code}[1]{\texttt{\small #1}}
\newcommand{\axiomref}[1]{\textbf{A#1}}
\newcommand{\msg}{\textsf{MSG}}
\newcommand{\ssg}{\textsf{SSG}}
\newcommand{\abb}{\textsf{ABB}}
\newcommand{\msgop}{\mathrm{MSG}}
\newcommand{\closer}[1]{C_R(#1)}
\newcommand{\Real}{\mathbb{R}}
\newcommand{\Realp}{\Real_{\geq 0}}
\newcommand{\inset}[1]{\mathrm{in}(#1)}
\newcommand{\outset}[1]{\mathrm{out}(#1)}
\newcommand{\bw}{\boldsymbol{w}}
\newcommand{\bc}{\boldsymbol{c}}

\setlist[itemize]{leftmargin=1.4em,itemsep=2pt,topsep=3pt}
\setlist[enumerate]{leftmargin=1.6em,itemsep=2pt,topsep=3pt}

\title{Extending the P-graph Formalism\\[4pt]
       \large\normalfont
       multi-objective $\abb$, signed-cycle correspondence,\\
       and open generalisations}
\author{For the $\HyMeKo${} project,\\
        in dialogue with the Friedler P-graph community}
\date{2026-05-19}

\begin{document}
\maketitle

\begin{abstract}\small
The Friedler P-graph formalism (Friedler, Tarj\'an, Huang, Fan 1992)
solves the cost-minimising chemical-process-synthesis problem via
two structural bounds (inclusion and reachability) over a bipartite
material/operating-unit hypergraph. We give a self-contained
restatement that exposes the abstract structure independent of the
chemical-process domain, then extend the formalism along three axes:
(i) multi-objective $\abb$ via positive-weighted scalarisation, with
an admissibility-preservation theorem; (ii) a structural
correspondence between Friedler $\abb$ and the signed-cycle
top-$K$ enumeration $\abb$ developed in this project; and
(iii) four open generalisations --- hierarchical, time-indexed,
stochastic, and categorical --- with axiom sketches and the natural
next theorems for each. Every result that has shipping code is
called out; everything else is sketched at the formal level
sufficient for the next iteration to be written.
\end{abstract}

\tableofcontents

\vspace{6pt}\hrule\vspace{6pt}

% ===================================================================
\section{Preliminaries and notation}
\label{sec:prelim}

\begin{definition}[P-graph]
\label{def:pgraph}
A \emph{P-graph} is a 6-tuple
$
\mathcal{P} \;=\; (M,\, O,\, \inset{\cdot},\, \outset{\cdot},\, R,\, P)
$
where:
\begin{itemize}
\item $M$ is a finite set of \emph{materials};
\item $O$ is a finite set of \emph{operating units}, disjoint from $M$;
\item $\inset{\cdot} : O \to 2^M$ assigns to each unit its
      \emph{input materials} (consumed);
\item $\outset{\cdot} : O \to 2^M$ assigns to each unit its
      \emph{output materials} (produced);
\item $R \subseteq M$ is the \emph{raw-material} set (feedstocks);
\item $P \subseteq M$ is the \emph{product-demand} set
      (required outputs).
\end{itemize}
The directed bipartite incidence graph is
$
E_{\text{dir}} \;=\;
\{ (m, u) \;:\; u \in O,\; m \in \inset{u}\} \;\cup\;
\{ (u, m) \;:\; u \in O,\; m \in \outset{u}\}.
$
\end{definition}

\begin{definition}[Producibility closure]
\label{def:closure}
Given $O' \subseteq O$, the \emph{producibility closure} of the raw
set under $O'$ is the least fixed point
\[
\closer{O'} \;=\; \mathrm{lfp}\,\Bigl[\,X \;\mapsto\; R \,\cup\,
                \{\,m \in M \;:\; \exists u \in O',\; \inset{u} \subseteq X,\; m \in \outset{u}\,\}\,\Bigr].
\]
Operationally, $\closer{O'}$ is built by iteratively adding the
outputs of every unit whose inputs are already producible from $R$.
The iteration converges in $O(|O|)$ rounds.
\end{definition}

\begin{definition}[Feasibility]
\label{def:feasible}
A subset $O' \subseteq O$ is a \emph{solution structure} if:
\begin{enumerate}[label=(\alph*)]
\item \textbf{Input consistency:} for every $u \in O'$,
      $\inset{u} \subseteq \closer{O'}$.
\item \textbf{Demand:} $P \subseteq \closer{O'}$.
\item \textbf{Strict no-excess (optional):} every produced material
      not in $R \cup P$ is consumed by some $u \in O'$.
\end{enumerate}
The set of all solution structures is denoted $\mathfrak{S}(\mathcal{P})$.
\end{definition}

\begin{definition}[Cost interface]
\label{def:cost}
A \emph{cost interface} on $\mathcal{P}$ is a non-negative function
$c : O \to \Realp$. The \emph{total cost} of $O' \subseteq O$ is
$
\mathrm{tc}_c(O') = \sum_{u \in O'} c(u).
$
The \emph{cost-minimising synthesis problem} is
\begin{equation}
\label{eq:scalar-objective}
\min_{O' \in \mathfrak{S}(\mathcal{P})}\; \mathrm{tc}_c(O').
\end{equation}
\end{definition}

% ===================================================================
\section{The classical Friedler axioms}
\label{sec:classical}

We restate the five axioms with explicit logical content; the
1992 paper's prose is slightly compressed.

\begin{axiom}[Product-membership]
\label{ax:product}
Every required product $p \in P$ is a material:
$P \subseteq M$.
\end{axiom}

\begin{axiom}[Raw-reachability]
\label{ax:raw}
Every material in $M$ is either raw or producible by some unit
chain from raws:
$M \subseteq \closer{O}.$
\end{axiom}

\begin{axiom}[Unit-catalogue closure]
\label{ax:unit}
Every $u \in O$ has $\inset{u} \cup \outset{u} \subseteq M$ and
$\inset{u} \cap \outset{u} = \emptyset$.
\end{axiom}

\begin{axiom}[Connectedness]
\label{ax:conn}
For every $u \in O$, both $|\inset{u}| \geq 1$ and
$|\outset{u}| \geq 1$ (no orphan units).
\end{axiom}

\begin{axiom}[No-cycle on materials]
\label{ax:acycle}
The directed bipartite incidence graph $E_{\text{dir}}$ has no
directed cycle that visits a material twice. (Equivalently,
$E_{\text{dir}}$ is a DAG when contracted on materials.)
\end{axiom}

\begin{remark}
\axiomref{1} and \axiomref{2} are well-formedness conditions on
the problem statement; \axiomref{3}--\axiomref{5} are
\emph{structural feasibility} preconditions that any P-graph
admitted to $\msg$ must satisfy. Stage P-mo's lowering enforces
\axiomref{3}--\axiomref{4} at parse time; \axiomref{5} is enforced
by the bipartite-schema invariant in $\HyMeKo$.
\end{remark}

% ===================================================================
\section{MSG, SSG, ABB as operators on the subset lattice}
\label{sec:operators}

\begin{definition}[Maximal structure $\msgop(\mathcal{P})$]
\label{def:msg}
The \emph{maximal structure} is
$
\msgop(\mathcal{P}) \;=\;
\bigcup_{O' \in \mathfrak{S}(\mathcal{P})} O'.
$
Equivalently (Friedler 1992), $\msgop$ is the largest $O'$ that
survives the fixed-point iteration of forward-feasibility
(drop $u$ whose inputs are not in $\closer{O'}$) and
backward-utility (drop $u$ whose outputs are not in the union of
inputs of remaining units $\cup P$) trims.
\end{definition}

\begin{definition}[$\ssg$]
$\ssg(\mathcal{P})$ enumerates every $O' \subseteq \msgop(\mathcal{P})$
that satisfies feasibility (Def.~\ref{def:feasible}). For
$|\msgop(\mathcal{P})| > 30$, $\ssg$ is impractical and $\abb$
is the right tool.
\end{definition}

\begin{definition}[$\abb$]
$\abb(\mathcal{P}, c)$ returns
$\arg\min_{O' \in \mathfrak{S}(\mathcal{P})} \mathrm{tc}_c(O')$
via depth-first include/exclude branching on a fixed ordering of
$\msgop(\mathcal{P})$, with two bounds:

\begin{enumerate}[label=\textbf{B\arabic*}]
\item \label{B:incl}
\textbf{Inclusion bound.} Let $s.\mathrm{cost}$ be the partial
sum at the current node and $s.\mathrm{best}$ the cost of the
current incumbent. If
$s.\mathrm{cost} \geq s.\mathrm{best}$, prune.
\item \label{B:reach}
\textbf{Reachability bound.} Let
$\widehat{O} = \mathrm{included} \cup \mathrm{undecided}$.
If $\closer{\widehat{O}} \not\supseteq P$, prune.
\end{enumerate}

Branching is include-before-exclude so the inclusion bound has
an incumbent to compare against as early as possible.
\end{definition}

\begin{proposition}[$\abb$ correctness]
\label{prop:abb-correct}
If both \textbf{\ref{B:incl}} and \textbf{\ref{B:reach}} are admissible
(prune only nodes whose subtree cannot improve the incumbent), then
$\abb(\mathcal{P}, c)$ returns an optimal feasible solution
structure, or reports infeasibility if $\mathfrak{S}(\mathcal{P}) = \emptyset$.
\end{proposition}

\begin{proof}[Proof sketch]
\textbf{\ref{B:incl}} is admissible because every inclusion adds
a non-negative cost, so the partial cost lower-bounds the
descendant's full cost. \textbf{\ref{B:reach}} is admissible
because $\closer{\widehat{O}} \supseteq \closer{O'}$ for every
$O' \subseteq \widehat{O}$ (monotonicity of $C_R$ in its argument),
so if $\closer{\widehat{O}} \not\supseteq P$, no descendant subset
can be feasible. Since both bounds preserve every optimum, the
DFS visits at least one optimum and returns the minimum.
\end{proof}

% ===================================================================
\section{Multi-objective extension: positive-weighted scalarisation}
\label{sec:multi-objective}

We now lift the cost interface from a scalar function
$c : O \to \Realp$ to a $D$-dimensional vector function
$\bc : O \to \Realp^D$ across named dimensions
(e.g.\ CAPEX, OPEX, CO$_2$, H$_2$O).

\begin{definition}[Multi-objective cost interface]
\label{def:multi-cost}
A \emph{$D$-dimensional cost interface} on $\mathcal{P}$ is a
function $\bc : O \to \Realp^D$, written
$\bc(u) = (c_1(u), \ldots, c_D(u))$. Given a user-supplied
\emph{weight vector} $\bw \in \Realp^D$, the
\emph{weighted-sum cost} of $u \in O$ is
$
c_{\bw}(u) \;=\; \langle \bw, \bc(u)\rangle
\;=\;\sum_{d=1}^{D} w_d \cdot c_d(u),
$
and the weighted total cost of $O' \subseteq O$ is
$
\mathrm{tc}_{\bw}(O') \;=\; \sum_{u \in O'} c_{\bw}(u).
$
\end{definition}

\begin{theorem}[Admissibility preservation under positive-weighted scalarisation]
\label{thm:admiss}
For every $\bw \in \Realp^D$ and every multi-objective cost
interface $\bc : O \to \Realp^D$, the $\abb$ inclusion bound
\textbf{\ref{B:incl}} restated as
\[
\text{prune when } \sum_{u \in \mathrm{included}} c_{\bw}(u)
\;\geq\; s.\mathrm{best}_{\bw}
\]
is admissible, and Proposition~\ref{prop:abb-correct} extends:
$\abb(\mathcal{P}, \bc, \bw)$ returns
$\arg\min_{O' \in \mathfrak{S}(\mathcal{P})} \mathrm{tc}_{\bw}(O')$.
\end{theorem}

\begin{proof}
Admissibility of \textbf{\ref{B:incl}} requires that the partial
weighted cost be a lower bound on every descendant's full weighted
cost. Since $w_d \geq 0$ for every $d$ and $c_d(u) \geq 0$ for every
$d, u$, every inclusion-step term
$c_{\bw}(u) = \sum_d w_d c_d(u)$ is non-negative. The partial sum
is therefore monotone-non-decreasing along the include-edge of the
search tree, so the partial weighted cost lower-bounds the
descendant's full weighted cost. The reachability bound
\textbf{\ref{B:reach}} is structural and does not depend on the
cost interface; it is unaffected. The rest of
Prop.~\ref{prop:abb-correct} applies verbatim.
\end{proof}

\begin{remark}[Why non-negativity matters]
Theorem~\ref{thm:admiss} fails for general $\bw \in \Real^D$. If
some $w_d < 0$, including a unit can decrease the partial cost,
and the inclusion bound no longer holds. A \emph{benefit term}
(negative weight, e.g.\ revenue) can be made admissible by
treating it separately --- subtract the maximum possible benefit
of the remaining subtree from the bound, see Section~\ref{sec:open}.
\end{remark}

\begin{corollary}[Implementation status]
\label{cor:impl}
The reference implementation in
\code{hymeko\_pgraph::abb} (commit Stage P-mo, 2026-05-19)
realises Theorem~\ref{thm:admiss} verbatim. The new
$\abb$Options field
\code{cost\_weights: Option<Vec<f64>>} dispatches between the
scalar path (Friedler 1992, byte-identical) and the
weighted-sum path. The methanol-synthesis worked example
(\code{data/pgraph/methanol\_synthesis.hymeko}) demonstrates
four structurally distinct optima from five weight regimes,
including a non-obvious hybrid topology that no single-criterion
regime selects.
\end{corollary}

% ===================================================================
\section{The signed-cycle correspondence}
\label{sec:cycle-corr}

Beyond chemical-process synthesis, this project has independently
applied $\abb$-style branch-and-bound to a structurally different
problem: top-$K$ enumeration of signed cycles in a signed graph
for graph-neural-network training. We now show these two
$\abb$ instances are not analogues but \emph{the same abstract algorithm}.

\begin{definition}[Top-$K$ signed-cycle enumeration]
\label{def:cycle-abb}
Given a signed graph $G = (V, E, \sigma : E \to \{-1, +1\})$, a
cycle length $k$, a target heap size $K$, and a per-cycle score
function $s : \mathrm{Cyc}_k(G) \to \Real$, the
\emph{top-$K$ enumeration problem} returns the $K$ closed
cycles of length $k$ with the largest scores. ABB uses a
\emph{bounded scorer trait} $(s, \overline{s})$, where
$\overline{s}(n_{\text{neg}}, k_{\text{rem}}, k)$ is an upper bound
on the score of any closed extension of a partial path with
$n_{\text{neg}}$ negative edges so far and $k_{\text{rem}}$
remaining edges to close, and prunes when
$\overline{s} \leq \mathrm{heap.peek\_min}$.
\end{definition}

\begin{theorem}[Cycle-PSE correspondence]
\label{thm:cycle-pse}
The cycle-enumeration $\abb$ (Def.~\ref{def:cycle-abb}) and the
P-graph $\abb$ (Section~\ref{sec:operators}) are both instances of
the following abstract algorithm. Let
$\mathcal{T}$ be a finite search tree with leaves
$\mathcal{L}$; let $\mathrm{obj} : \mathcal{L} \to \Real$ be an
objective; let
$\overline{\mathrm{obj}} : \mathrm{nodes}(\mathcal{T}) \to \Real$
be an admissible upper bound (or lower bound for minimisation) on
the optimum over each subtree. Then \textbf{Algorithm~A1}
\textsf{[DFS-with-bound]}
visits at least one optimum and returns it in time
$O(|\mathcal{T}_{\overline{\mathrm{obj}}\text{-feasible}}|)$, where
$|\mathcal{T}_{\overline{\mathrm{obj}}\text{-feasible}}|$ counts the
nodes whose bound does not allow pruning.
\end{theorem}

\begin{proof}[Proof sketch]
In the cycle case, $\mathcal{L}$ is the set of all closed cycles
of length $k$ from a chosen start vertex, $\mathrm{obj}$ is the
per-cycle score, and the bound is the BoundedScorer's
\code{upper\_bound} method.
In the PSE case, $\mathcal{L}$ is the set of feasible solution
structures, $\mathrm{obj}$ is the (possibly weighted) cost, and the
bound is the partial-cost-plus-zero-future = partial-cost itself
(plus the structural reachability gate). Both fit Algorithm~A1
with $\mathcal{T}$ being a binary include/exclude tree (PSE) or a
$|V|$-way edge-extend tree (cycles).
\end{proof}

\begin{remark}
The shared abstract algorithm explains why the implementation
patterns parallel each other so closely. The \HyMeKo{} codebase
duplicates structural code across the two $\abb$ instances; a
plausible refactor (not in scope here) would extract a
\code{trait BoundedSearch} that both inhabit.
\end{remark}

% ===================================================================
\section{Open generalisations}
\label{sec:open}

We sketch four extensions, each with the natural next axiom or
theorem and a pointer to the next implementation step.

\subsection{Pareto-front enumeration}
\label{ssec:pareto}

For a multi-objective P-graph $(\mathcal{P}, \bc, D)$, the
\emph{Pareto front} $\mathcal{F} \subseteq \mathfrak{S}(\mathcal{P})$ is the
set of feasible structures $O'$ such that no other feasible
$O''$ satisfies
$c_d(O'') \leq c_d(O')$ for all $d$
with strict inequality for at least one. Theorem~\ref{thm:admiss}
gives \emph{one} Pareto-optimal point per $\bw$; sweeping $\bw$
over the unit simplex enumerates the whole front.

\begin{proposition}[Simplex-grid completeness for $D \leq 3$]
\label{prop:simplex}
For $D \leq 3$ and bounded integer cost vectors, the Pareto front
$\mathcal{F}$ is finite, and a simplex-grid sweep of $\bw$ at
resolution $1/N$ recovers every Pareto-optimal structure for
sufficiently large $N$ (depending on the gaps between adjacent
Pareto points).
\end{proposition}

\begin{proof}[Proof sketch]
The Pareto front of finitely many $\Realp^D$ points is itself
finite. For each Pareto point $O^\star$, the cone of weights
$\bw$ for which $O^\star$ is the weighted optimum is open and
non-empty (otherwise $O^\star$ is dominated). A fine-enough simplex
grid intersects every such cone.
\end{proof}

\textbf{Implementation step:} $\sim 100$ LOC, builds on the existing
SSG enumeration loop. For $D \geq 4$ the simplex grid is
exponential in $D$ and sampling becomes more efficient than grid;
adaptive ($\epsilon$-constraint or weighted-$L^2$-Tchebycheff)
methods take over.

\subsection{Hierarchical P-graphs}

A natural generalisation models each operating unit as itself a
P-graph (a sub-plant). Formally:

\begin{definition}[Hierarchical P-graph]
\label{def:hier}
A \emph{hierarchical P-graph} is a tuple
$\mathcal{H} = (M, O, \inset{\cdot}, \outset{\cdot}, R, P, \mathrm{sub})$
where $\mathrm{sub}: O \to \mathfrak{P}_{\text{hier}}$ is a
(possibly partial) function assigning to each unit a sub-plant
(itself a hierarchical P-graph). The flat unrolling of $\mathcal{H}$
substitutes each $u \in O$ with $\mathrm{sub}(u)$, plumbing the
boundary materials through.
\end{definition}

\begin{proposition}[Hierarchical $\abb$]
\label{prop:hier-abb}
$\abb$ on a hierarchical P-graph $\mathcal{H}$ reduces to nested
$\abb$ calls: solve each sub-plant for the cost it contributes to
its parent unit, then run parent-level $\abb$ with the cached
sub-plant costs. The reachability bound becomes a
\emph{nested closure}: a parent unit is reachable iff its
sub-plant has a feasible structure under the parent's input
materials.
\end{proposition}

\textbf{Implementation step:} sub-plant caching plus a recursive
$\abb$ entry. The bound for the nested closure requires care;
the natural extension is to compute the sub-plant's minimum cost
over its input/output interface as the parent-level cost.

\subsection{Time-indexed P-graphs}

Real plants operate in time. We add an explicit time index.

\begin{definition}[Time-indexed P-graph]
\label{def:time}
A \emph{time-indexed P-graph} over a finite horizon
$T = \{0, 1, \ldots, H\}$ is a tuple
$(\mathcal{P}, \tau_{\mathrm{in}}, \tau_{\mathrm{out}}, S)$
where:
\begin{itemize}
\item $\tau_{\mathrm{in}}, \tau_{\mathrm{out}} : O \to \mathbb{Z}_{\geq 0}$
      are per-unit consumption and production latencies;
\item $S \subseteq M$ is a set of \emph{storable} materials
      (others must be consumed in the same time-step as produced).
\end{itemize}
A \emph{time-indexed solution structure} is a function
$\Sigma : T \to 2^O$ assigning to each time-step the set of units
active at that step, subject to a time-extended feasibility:
inputs at time $t$ must be available from raws or from outputs
produced no later than $t - \tau_{\mathrm{in}}(u)$, and outputs
of non-storable materials must be consumed at exactly
$t + \tau_{\mathrm{out}}(u)$.
\end{definition}

\begin{axiom}[Time-monotonicity / A6]
\label{ax:time}
For every storable material $m \in S$ and every time-step $t$,
the available inventory of $m$ at $t$ is a non-decreasing function
of the production schedule up to $t$. (Production is irreversible;
no negative inventory.)
\end{axiom}

\textbf{Implementation step:} the rolling-horizon $\abb$ takes
$\Sigma : T \to 2^O$ as the lattice, with the reachability bound
generalised to per-time-step closure under inventory dynamics.
This is non-trivial work; the natural prototype is $H \leq 24$
hourly steps on the methanol example.

\subsection{Stochastic P-graphs}

Real plants face uncertain feedstock costs, demand, and
equipment availability. We treat the cost interface as random.

\begin{definition}[Stochastic P-graph]
\label{def:stoch}
A \emph{stochastic P-graph} is a tuple
$(\mathcal{P}, \mathbb{P})$ where $\mathbb{P}$ is a probability
measure on a sample space $\Omega$ and the cost interface is a
random function
$c : O \times \Omega \to \Realp$. The decision variable is the
structure $O'$; the realisation $\omega \in \Omega$ is drawn
\emph{after} the decision. The \emph{risk-neutral} objective is
$
\min_{O'} \mathbb{E}_\omega[\mathrm{tc}_c(O', \omega)].
$
The \emph{conditional-value-at-risk} (CVaR) objective at level
$\alpha \in (0, 1]$ is
$
\min_{O'} \mathrm{CVaR}_\alpha[\mathrm{tc}_c(O', \omega)].
$
\end{definition}

\begin{axiom}[Risk-coherence / A7]
\label{ax:risk}
The risk functional $\rho : L^1(\Omega, \mathbb{P}) \to \Real$
employed in the stochastic objective is \emph{coherent}: monotone,
sub-additive, positively homogeneous, and translation-invariant
(Artzner et al.\ 1999). $\mathbb{E}$ and $\mathrm{CVaR}_\alpha$ both
satisfy A7.
\end{axiom}

\begin{proposition}[Sample-average approximation]
\label{prop:saa}
For coherent $\rho$ and a sample $\omega_1, \ldots, \omega_N$ drawn
i.i.d.\ from $\mathbb{P}$, the sample-average $\abb$
\[
\min_{O'} \rho\bigl(\{\mathrm{tc}_c(O', \omega_n)\}_{n=1}^N\bigr)
\]
satisfies an admissibility-preservation theorem analogous to
Theorem~\ref{thm:admiss}, because $\rho$ being monotone and
translation-invariant preserves the cost-monotonicity of the
inclusion bound.
\end{proposition}

\textbf{Implementation step:} per-realisation cost vectors as
multi-objective cost vectors of size $D = N$; the weight vector
encodes the risk functional. CVaR$_\alpha$ corresponds to taking
the largest $\lceil N \alpha\rceil$ realisations' average; this can
be re-cast as a weighted-sum with a structural weight constraint,
fitting Theorem~\ref{thm:admiss} after a permutation step.

\subsection{Categorical formulation}

The bipartite structure of a P-graph and the compositional nature
of nested sub-plants suggest a categorical view. Sketch only:

\begin{definition}[$\mathbf{PGraph}$]
Let $\mathbf{PGraph}$ be the category whose objects are finite
sets of materials (with raw/product designations) and whose
morphisms $\phi : M_R \to M_P$ are P-graphs from $R = M_R$ to
$P = M_P$. Composition $\phi \circ \psi$ glues two P-graphs by
identifying the product set of $\psi$ with the raw set of $\phi$.
The identity morphism is the empty P-graph
($O = \emptyset$, $R = P$).
\end{definition}

\begin{proposition}[$\mathbf{PGraph}$ is symmetric monoidal]
The disjoint union of P-graphs equips $\mathbf{PGraph}$ with a
symmetric monoidal structure
$\otimes : \mathbf{PGraph} \times \mathbf{PGraph} \to \mathbf{PGraph}$,
with unit $I = (\emptyset, \emptyset, \emptyset, \emptyset)$. Recycle
loops correspond to the trace operator
$\mathrm{Tr}^X : \mathbf{PGraph}(A \otimes X, B \otimes X) \to \mathbf{PGraph}(A, B)$,
making $\mathbf{PGraph}$ a traced symmetric monoidal category.
\end{proposition}

\begin{remark}
The categorical view connects P-graphs to other formalisms in the
SMC ecosystem (Petri nets, signal-flow graphs, string diagrams).
The trace structure formalises recycle loops as morphism feedback;
the multi-objective cost functor sends $\mathbf{PGraph}$ to
$\mathbf{Lat}^D$ (the category of $D$-component lattices), with
$\abb$ as the value-recovery functor's adjoint. This connection
is exploratory and not yet implemented.
\end{remark}

\textbf{Implementation step:} none in the current codebase. The
categorical framing is a paper-only generalisation that would
inform a future redesign of the IR shared between $\abb$ instances.

% ===================================================================
\section{Implementation status and roadmap}
\label{sec:status}

\begin{center}\small
\begin{tabular}{lccc}
\toprule
Extension & Section & Status & Estimated LOC \\
\midrule
Scalar $\abb$ (Friedler 1992) & \ref{sec:operators} & shipped & --- \\
Multi-objective $\abb$ (Thm.~\ref{thm:admiss})       & \ref{sec:multi-objective} & shipped Stage P-mo 2026-05-19 & $\sim 120$ \\
Cycle-PSE correspondence (Thm.~\ref{thm:cycle-pse})  & \ref{sec:cycle-corr}      & shipped (both ends) & --- \\
Pareto-front enumeration (Prop.~\ref{prop:simplex})  & \ref{ssec:pareto}         & queued & $\sim 100$ \\
Hierarchical P-graphs (Prop.~\ref{prop:hier-abb})    & \ref{sec:open}            & queued & $\sim 300$ \\
Time-indexed P-graphs (\axiomref{6})                 & \ref{sec:open}            & sketched & $\sim 800$ \\
Stochastic P-graphs (\axiomref{7}, Prop.~\ref{prop:saa}) & \ref{sec:open}        & sketched & $\sim 400$ \\
Categorical $\mathbf{PGraph}$                        & \ref{sec:open}            & exploratory & --- \\
\bottomrule
\end{tabular}
\end{center}

\paragraph{What the formalism is actually used for in this project.}
\begin{enumerate}
\item Process synthesis on the canonical HDA example (test fixture).
\item Multi-objective methanol synthesis (\emph{this report's worked example}).
\item Top-$K$ signed-cycle enumeration in HSiKAN (graph-neural-net
      training); shares the abstract Algorithm~A1.
\item Graph-property-conditional axiom selection for signed-link
      prediction; the regime-conditional rule is structurally an
      $\abb$-style bound choice (which axiom to enforce given the
      graph's balance ratio).
\end{enumerate}

% ===================================================================
\section{Bottom line}
\label{sec:bottom}

The Friedler 1992 P-graph formalism is robust enough to bear
serious extensions. We've shipped two: the multi-objective ABB
with admissibility preservation (Theorem~\ref{thm:admiss}) and the
structural correspondence to signed-cycle enumeration
(Theorem~\ref{thm:cycle-pse}). Four further generalisations
(Pareto-front, hierarchical, time-indexed, stochastic) are
formally sketched with explicit axioms or theorems and a clear
implementation step for each. A categorical formulation hints at
a deeper unification but is exploratory.

The formalism's value for the PSE community goes beyond its
original chemical-process-synthesis domain: any combinatorial
decision problem where (a) the search space is the subset lattice
of a finite catalogue, (b) feasibility is closed under reachability
of a target set, and (c) the objective is non-negatively
weighted-additive over the catalogue, is an instance of
P-graph $\abb$. The cycle-enumeration application is a proof of
concept that this generality is operational. The next big
opportunities --- time-indexed plant scheduling, stochastic
feedstock procurement, hierarchical sub-plant nesting --- are all
within the formal envelope established here.

\vspace{8pt}\hrule\vspace{3pt}
\noindent\footnotesize
\textbf{References to the implementation.} Plan
\code{docs/plans/2026-05-19-pgraph-multi-objective/};
implementation in
\code{hymeko\_pgraph/src/\{lowering,abb,dump\}.rs};
9 new tests in
\code{hymeko\_pgraph/tests/multi\_objective.rs};
worked example in
\code{data/pgraph/methanol\_synthesis.hymeko};
companion 5-page demo PDF in
\code{reports/2026-05-19-pgraph-multi-objective.pdf}. The
top-$K$ signed-cycle $\abb$ is implemented in
\code{hymeko\_graph/src/topk\_cycles.rs}; see
\code{reports/2026-05-10-abb-global-topk.md} for the
$25.06\times$ Epinions benchmark.

\medskip
\noindent
This document is intended as input for an upcoming paper
``P-graph methodology beyond PSE: axiom-feasibility for
sustainable combinatorial optimisation'' (working title) to be
discussed with the Friedler community, including Jean Pimentel.
Feedback on the axiomatic statements, the categorical framing,
and the priority among the open generalisations is welcome.

\end{document}
